\documentclass[]{ling-200b-2}
\usepackage[utf8]{inputenc}
\usepackage[greek.ancient,english]{babel}
\usepackage{alphabeta}
\usepackage{tipa}
\usepackage{bold-extra}
\usepackage{ulem}

\usepackage[letterpaper, margin = 1in]{geometry}
\usepackage{fancyhdr}
\usepackage{titling}
\usepackage{titlesec}
\usepackage{hyperref}
\usepackage{lastpage}
%\usepackage{lipsum}
\usepackage{amssymb}
\usepackage{setspace}
\usepackage{stmaryrd}
\usepackage{rotating}
\usepackage{lscape}
\usepackage{amsmath}
\usepackage[math-style=ISO]{unicode-math}

\usepackage{hanging}
\usepackage{tabularx}
\usepackage{multirow}
\usepackage{array}
\renewcommand\tabularxcolumn[1]{m{#1}}
\usepackage{arydshln}
\usepackage[table]{xcolor}

\usepackage{tikz}
\usepackage{tikz-qtree}
\usepgflibrary{arrows}
\usetikzlibrary{positioning,shapes.multipart}
\usetikzlibrary{tikzmark}
\usetikzlibrary{fit}
%\usepackage{tree-dvips}

\usepackage[backend=biber, style=apa]{biblatex}
\setlength{\bibhang}{0.5in}
\addbibresource{205.bib}
\usepackage{csquotes}
\DeclareLanguageMapping{english}{english-apa}

%\usepackage{linguex}
%\usepackage{gb4e}

\title{Attic Greek Aorist/Future Passives}
\author{Ian Thomas White}
\date{02/21/2021}

\makeatletter
\renewcommand{\maketitle}
{
\begin{flushleft}
\textbf{{\large\@title}}
\newline
\@author
\end{flushleft}
}
\makeatother

\titleformat*{\section}{\medium\bfseries}
\titleformat*{\subsection}{\small\bfseries}

\newcommand*{\textsup}[1]{\ensuremath{^\textrm{{\scriptsize #1}}}}
\newcommand*{\textsub}[1]{\ensuremath{_\textrm{{\scriptsize #1}}}}
\newcommand*{\tinytextsup}[1]{\ensuremath{^\textrm{{\tiny #1}}}}
\newcommand*{\tinytextsub}[1]{\ensuremath{_\textrm{{\tiny #1}}}}

\pagestyle{fancy}
\fancyhf{}
\lhead{Ian Thomas White}
\chead{}
\rhead{\thepage \ of \pageref{LastPage}}

% \linespread{1.6}
\widowpenalty=100000
\clubpenalty=100000

\usepackage{enumitem}
\usepackage{verbatim}
\newcommand{\pres}[0]{Tns$_\textsc{pres}$}
\newcommand{\past}[0]{Tns$_\textsc{past}$}
\usetikzlibrary{decorations.pathreplacing}

\renewcommand{\labelenumi}{$\cdot$\Roman{enumi}$\cdot$}
\renewcommand{\labelenumii}{$\cdot$\Alph{enumii}$\cdot$}
\renewcommand{\labelenumiii}{$\cdot$\roman{enumiii}$\cdot$}

\usepackage{graphicx}
\graphicspath{ {./images/} }

\begin{document}

\thispagestyle{empty}
\singlespacing
\maketitle

\section{Introduction}

Despite the robust literature on Modern Greek and Classical Latin, even in the DM world, relatively little formal attention has been given to Classical Greek, at least in the generative tradition. The current paper attempts to address the gap by analyzing a rather puzzling, and puzzled over, paradigm from Attic Greek: the aorist/future passive. Namely, I provide a DM analysis of curious deponent phenomena in the regular Greek (thematic) verb: the Greek verb has distinct inflectional endings for active versus middle/passive forms, but the aorist passive requires the active endings. Furthermore, the future passive, which shares a suffix with the aorist passive, requires the otherwise expected middle/passive endings. Given certain assumptions about the structure of the Greek and the features obtaining on the Greek voice head, I show that these facts can be accounted under DM using spans, VI underspecification, and redundancy rules.

The organization of the paper is as follows. Section 2 gives the relevant paradigms for the puzzle. Section 3 discuss the structure and features of the Attic verb that are necessary to account for the data. Section 4 presents the proposed solution. Finally, section 5 offers some concluding remarks.


\section{Paradigms}

Attic Greek, as traditionally described, has four tense-aspect categories (present, future, aorist, perfect). Each tense-aspect is broken into two or three voice categories: active versus middle/passive in the present and perfect and active versus middle versus passive in the future and aorist. Finally, there are two sets inflectional endings that manifest, in some way, voice: active and middle/passive. The inflectional endings always come at the right edge of the verb. Markings for other features are less clear and are discussed below. I give the full paradigm for the finite forms of the regular thematic verb in (1) and (2).\footnote{Thematic verbs have a theme vowel preceding the inflectional endings, contrastive to athematic verbs.}

\noindent \hfill \par
\ex[numoffset=\leftmargin] \textit{3rd sing. across tense/mood/aspect for} λύειν  `release/kill'
\xe
\scriptsize
\noindent \hfill \par
\noindent\begin{tabularx}{0.95\textwidth}{
l
l
| >{\centering\arraybackslash}X
>{\centering\arraybackslash}X
>{\centering\arraybackslash}X
>{\centering\arraybackslash}X
>{\centering\arraybackslash}X
}
& & indic. 1 & indic. 2 & subj. & opt. & imper. \\
\hline
\multirow{2}{3em}{present} & act. & λύ-ε-ι & ἔ-λυ-ε-$\varnothing$ & λύ-ῃ & λύ-οι-$\varnothing$ & λυ-έ-τω \\
 & mid.-pass. & λύ-ε-ται & ἐ-λύ-ε-το & λύ-η-ται & λύ-οι-το & λυ-έ-σθω \\
\hline
\multirow{3}{3em}{future} & act. & λύ-σ-ε-ι & - & - & λύ-σ-οι-$\varnothing$ & - \\
 & mid. & λύ-σ-ε-ται & - & - & λύ-σ-οι-το & - \\
 & \color{red} pass. & \color{red} λύ-θή-σ-ε-ται & - & - & \color{red} λυ-θή-σ-οι-το & - \\
\hline
\multirow{3}{3em}{aorist} & act. & - & ἔ-λυ-σ-ε-$\varnothing$ & λύ-σ-ῃ & λύ-σ-αι-$\varnothing$ & λυ-σ-ά-τω \\
 & mid. & - & ἐ-λύ-σ-α-το & λύ-σ-η-ται & λύ-σ-αι-το & λυ-σ-ά-σθω \\
 & \color{red} pass. & - & \color{red} ἐ-λύ-θη-$\varnothing$ & \color{red} λυ-θῇ & \color{red} λυ-θ-εί-η & \color{red} λυ-θή-τω \\
\hline
\multirow{2}{3em}{perfect} & act. & λέ-λυ-κ-ε-$\varnothing$ & ἐ-λε-λύ-κ-ε-ι & λε-λύ-κ-ῃ & λε-λύ-κ-οι-$\varnothing$ & λε-λυ-κ-έ-τω \\
 & mid.-pass. & λέ-λυ-ται & ἐ-λέ-λυ-το & λε-λυ-μένος ᾖ & λε-λυ-μένος εἴη & λε-λύ-σθω \\
 \hline
 fut. perf. & mid.-pass. & λε-λύ-σ-ε-ται & - & - & ? & ? \\

\end{tabularx}

\normalsize
\noindent \hfill \par
\noindent \hfill \par
\ex[numoffset=\leftmargin] \textit{3rd sing. across tense/mood/aspect for} λύειν  `release/kill' (\textit{IPA})\footnote{Note that Greek ε is IPA [e], and Greek η is IPA [ε:]. The stress and breathing marks in the Greek are irrelevant here.}
\xe
\scriptsize
\noindent \hfill \par
\noindent\begin{tabularx}{\textwidth}{
l
l
| >{\centering\arraybackslash}X
>{\centering\arraybackslash}X
>{\centering\arraybackslash}X
>{\centering\arraybackslash}X
>{\centering\arraybackslash}X
}
& & indic. 1 & indic. 2 & subj. & opt. & imper. \\
\hline
\multirow{2}{3em}{present} & act. & ly-e-i & e-ly-e-$\varnothing$ & ly-ε:-i & ly-oi-$\varnothing$ & ly-e-t\textipa{O}: \\
 & mid.-pass. & ly-e-tai & e-ly-e-to & ly-ε:-tai & ly-oi-to & ly-e-st\textipa{\super h\nospace O}: \\
\hline
\multirow{3}{3em}{future} & act. & ly-s-e-i & - & - & ly-s-oi-$\varnothing$ & - \\
 & mid. & ly-s-e-tai & - & - & ly-s-oi-to & - \\
 & \color{red} pass. & \color{red} ly-t\textipa{\super h}ε:-s-e-tai & - & - & \color{red} ly-t\textipa{\super h}ε:-s-oi-to & - \\
\hline
\multirow{3}{3em}{aorist} & act. & - & e-ly-s-e-$\varnothing$ & ly-s-ε:-i & ly-s-ai-$\varnothing$ & ly-s-a-t\textipa{O}: \\
 & mid. & - & e-ly-s-a-to & ly-s-ε:-tai & ly-s-ai-to & ly-s-a-st\textipa{\super h\nospace O}:\\
 & \color{red} pass. & - & \color{red} e-ly-t\textipa{\super h}ε:-$\varnothing$ & \color{red} ly-t\textipa{\super h}ε:-i & \color{red} ly-t\textipa{\super h}-ei-ε: & \color{red} ly-t\textipa{\super h}-t\textipa{O}:  \\
\hline
\multirow{2}{3em}{perfect} & act. & le-ly-k-e-$\varnothing$ & e-le-ly-k-e-i & le-ly-k-ε:-i & le-ly-k-oi-$\varnothing$ & le-ly-k-e-t\textipa{O}: \\
 & mid.-pass. & le-ly-tai & e-le-ly-to & le-ly-menos ε:i & le-ly-menos eiε: & le-ly-st\textipa{\super h\nospace O}\\
 \hline
 fut. perf. & mid.-pass. & le-ly-s-e-tai & - & - & ? & ? \\

\end{tabularx}

\normalsize
\noindent \hfill \par
\noindent \hfill \par
Thus, there are four sets of agreement endings: active primary, middle-passive primary, active secondary, middle-passive secondary (where primary/secondary map onto nonpast/past), e.g., in (1,2) the present forms of the active indicative 1, active indicative 2,\footnote{The ``indicative 2'' is also called the ``imperfect''.} middle-passive indicative 1, middle-passive indicative 2, respectively. The aorist and future tense-aspects are unique in the paradigm in that they have separate forms for the passive and middle,  rather than a single middle-passive form. The passive as seen in (1,2) is marked by -θη- [-t\textsup{h}ε:-] affix on the root\footnote{ In some verbs the passive is marked only with -η-/[-ε:-]. I won't deal with those verbs here}. What is especially curious is that the aorist passive, unlike any other (middle-)passive in Greek, takes active (secondary) endings. On the other hand, the future passive, which, again, takes the same suffix as the aorist passive, takes the expected middle-passive (primary) endings. The relevant endings are marked in red in (3,4): the future middle and passive and the aorist passive all take middle-passive endings, while the aorist passive takes the active endings (most evident here in the 1st/2nd plural).\footnote{The endings themselves display a number of syncretisms that would be worth exploring, and they could be reasonably split into separate morphemes for person and number. I leave these tasks for future work.}

\noindent \hfill \par
\ex \textit{aorist/future active/middle/passive }
\xe

\scriptsize
\noindent \hfill \par
\begin{tabularx}{0.9\textwidth}{
l
| >{\centering\arraybackslash}X
 >{\centering\arraybackslash}X
 >{\centering\arraybackslash}X
| >{\centering\arraybackslash}X 
 >{\centering\arraybackslash}X
 >{\centering\arraybackslash}X
}
& \multicolumn{3}{c|}{future} & \multicolumn{3}{c}{aorist} \\
\cline{2-7}
& active & middle & passive & active & middle &  passive \\
\hline
1sg & λύ-σ-\color{red}ω & λύ-σ-ο-\color{red}μαι & λυ-θή-σ-ο-\color{red}μαι & ἔ-λυ-σ-α-\color{red}$\varnothing$ & ἐ-λυ-σ-ά-\color{red}μην & ἐ-λύ-θη-\color{red}ν \\
2sg & λύ-σ-ε-\color{red}ις & λύ-σ-ε-\color{red}ι/ῃ & λυ-θή-σ-ε-\color{red}ι/ῃ & ἐ-λυ-σ-\color{red}ω & ἐ-λύ-σ-\color{red}ω & ἐ-λύ-θη-\color{red}ς \\
3sg & λύ-σ-ε-\color{red}ι & λύ-σ-ε-\color{red}ται & λυ-θή-σ-ε-\color{red}ται & ἔ-λυ-σ-ε-\color{red}$\varnothing$ & ἐ-λύ-σ-α-\color{red}το & ἐ-λύ-θη-\color{red}$\varnothing$ \\
1pl & λύ-σ-ο-\color{red}μεν & λυ-σ-ό-\color{red}μεθα & λυ-θη-σ-ό-\color{red}μεθα & ἐ-λύ-σ-α-\color{red}μεν & ἐ-λυ-σ-ά-\color{red}μεθα & ἐ-λύ-θη-\color{red}μεν \\
2pl & λύ-σ-ε-\color{red}τε & λύ-σ-ε-\color{red}σθε & λυ-θή-σ-ε-\color{red}σθε & ἐ-λύ-σ-α-\color{red}τε & ἐ-λύ-σ-α-\color{red}σθε & ἐ-λύ-θη-\color{red}τε \\
3pl & λύ-σ-ου-\color{red}σι & λύ-σ-ο-\color{red}νται & λυ-θή-σ-ο-\color{red}νται & ἔ-λυ-σ-α-\color{red}ν & ἐ-λύ-σ-α-\color{red}ντο & ἐ-λύ-θη-\color{red}σαν \\
\end{tabularx}

\noindent \hfill \par
\noindent \hfill \par
\normalsize
\ex \textit{aorist/future active/middle/passive (IPA) }\footnote{ The present active 3rd person ending -ουσι [-o:si] is from -ο+νσι [-o+nsi], and the aorist middle 2nd person ending -σω [-s\textipa{O}:] is a contraction of -σα+ο [-sa+o].}  
\xe

\scriptsize
\noindent \hfill \par
\begin{tabularx}{0.9\textwidth}{
l
| >{\centering\arraybackslash}X
 >{\centering\arraybackslash}X
 >{\centering\arraybackslash}X
| >{\centering\arraybackslash}X
 >{\centering\arraybackslash}X
 >{\centering\arraybackslash}X
}
& \multicolumn{3}{c|}{future} & \multicolumn{3}{c}{aorist} \\
\cline{2-7}
& active & middle & passive & active & middle & passive \\
\hline
1sg & ly-s-\color{red}\textipa{O}: & ly-s-o-\color{red}mai & ly-t\textipa{\super h}ε:-s-o-\color{red}mai & e-ly-s-a-\color{red}$\varnothing$ & e-ly-s-a-\color{red}mε:n & e-ly-t\textipa{\super h}ε:-\color{red}n\\
2sg & ly-s-e-\color{red}is & ly-s-e/ε:-\color{red}i & ly-t\textipa{\super h}ε:-s-e/ε:-\color{red}i & e-ly-s-\color{red}\textipa{O}: & e-ly-s-\color{red}\textipa{O}: & e-ly-t\textipa{\super h}ε:-\color{red}s  \\
3sg & ly-s-e-\color{red}i& ly-s-e-\color{red}tai & ly-t\textipa{\super h}ε:-s-e-\color{red}tai & e-ly-s-e-\color{red}$\varnothing$ & e-ly-s-a-\color{red}to & e-ly-t\textipa{\super h}ε:-\color{red}$\varnothing$ \\
1pl & ly-s-o-\color{red}men & ly-s-o-\color{red}met\textipa{\super h}a & ly-t\textipa{\super h}ε:-s-o-\color{red}met\textipa{\super h}a & e-ly-s-a-\color{red}men & e-ly-s-a-\color{red}met\textipa{\super h}a & e-ly-t\textipa{\super h}ε:-\color{red}men \\
2pl & ly-s-e-\color{red}te & ly-s-e-\color{red}st\textipa{\super h}e & ly-t\textipa{\super h}-s-e-\color{red}t\textipa{\super h}e & e-ly-s-a-\color{red}te & e-ly-s-a-\color{red}st\textipa{\super h}e & e-ly-t\textipa{\super h}ε:-\color{red}te \\
3pl & ly-s-o:-\color{red}si & ly-s-o-\color{red}ntai & ly-t\textipa{\super h}ε:-s-o-\color{red}ntai & e-ly-s-a-\color{red}n & e-ly-s-a-\color{red}nto & e-ly-t\textipa{\super h}ε:-\color{red}san \\
\end{tabularx}

\normalsize
\section{The Greek Verb}

\subsection{Verb Structure}

Little formal work has been done on the structure of the Attic verb. On first blush it looks like the paradigm would include a number of instances of multiple exponence; e.g., the the aorist (active and middle) is marked by the augment prefix ε- [e-] in addition to the -σ- [-s-] affix, and the perfect (active) is marked by both prefixed reduplication and the -κ- [-k-] affix. The following claims have been made recently about which heads or features the various affixes of the Greek verb realize:

\noindent \hfill \par
\pex \textit{claims for realization of Greek affixes}
\a \textit{(temporal) augment} ε- [e-]: T[+past] (Schreiner 2021)
\a \textit{prefixed reduplication}: Asp[+perfect] (Schreiner 2021)
\a \textit{aorist} -σ- [-s-]: v/\underline{\hspace{1em}}Asp[+perfective] (Grestenberger 2015, 2019); Asp[+perfective] (Merchant 2015, Schreiner 2021)
\a \textit{future} -σ- [-s-]: T[+future] (Schreiner 2021); Mod(al)[+future] (Grestenberger 2014,2019)
\a \textit{perfect} -κ- [-k-]: Voice[+active] (Schreiner 2021)
\a \textit{passive} -θη- [-t\super hε:-]: Voice[Nonactive, passive] (Embick 1997); v/\underline{\hspace{1em}}[+perfective] (Grestenberger 2015, 2019)
\a \textit{theme vowels}: v (Grestenberger 2015, 2019); Mood (Schreiner 2021)
\a \textit{inflectional endings}: Agr (Embick 1997, Schreiner 2021); T+Agr (Grestenberger 2015, 2019)
\xe
\noindent \hfill \par

I will follow the general consensus (apart from Grestenberger 2015, 2019) such that the Greek verb has the basic structure in (6). This is a fairly common assumed syntactic structure cross-linguistically, as well. I will also assume, as is common  for verbal elements, that these heads are always present, even when phonologically null. I use V as a stand-in for any possible root-V-v complex. I will follow Embick (1997) in assuming that the passive suffix -θη- [-t\super hε:-] realizes Voice, as is intuitive, though I will use  slightly different features (see below). Otherwise, I'll follow Schreiner (2021) in assuming that the each affix realizes a distinct set of features, i.e., that there is no multiple exponence. In (7) I provide a breakdown of which head each affix realizes for the 3rd person indicative of λύω[ly\textipa{O}:]. I provide feature decompositions and explicit vocabulary items later.

\noindent \hfill \par
\ex \textit{general Greek verbal schema (post-HM)}

\small
\begin{forest}
[Agr
[Mood
[Tense
[Aspect
[Voice
[V]
[Voice]]
[Aspect]]
[Tense]]
[Mood]]
[Agr]]
\end{forest}
\xe

\noindent \hfill \par
\ex \textit{verbal affixes (indicative)}\footnote{ These are in their normal linearized order apart from the future σ/s which realizes T like the augment but shows up in a post-root position. Also, I've slotted the -σα- [-sa-] in the aorist as one Asp affix, rather than splitting it up into Asp and Mood affixes. This is the traditional conception of these affixes, but it is by no means clear. The way these vowels show up in the subjunctive and optative though could be evidence that it should be split up, but it makes the VIs below much more complicated, and nothing in the current analysis hinges on this decision.}
\xe

\noindent \hfill \par
\small
\begin{tabularx}{0.95\textwidth}{
>{\centering\arraybackslash}X
 >{\centering\arraybackslash}X
 >{\centering\arraybackslash}X
 >{\centering\arraybackslash}X
 >{\centering\arraybackslash}X
 >{\centering\arraybackslash}X
|c
} 
aug. & redup. & root & [suffix] & theme & infl. & parse \\
T & Asp & root & Voice & Mood & Agr & \\
\hline
 & & λυ/ly & & ε/e & ι/i & `present active' \\
 & & λυ/ly & & ε/e & ται/tai & `present middle-passive' \\
ε/e & & λυ/ly & & ε/e & $\varnothing$ & `imperfect active' \\
ε/e & & λυ/ly & & ε/e & το/to & `imperfect middle-passive' \\
σ/s & & λυ/ly & & ε/e & ι/i & `future active' \\
σ/s & & λυ/ly & & ε/e & ται/tai & `future middle' \\
 σ/s & & λυ/ly & θη/t\super hε: & ε/e & ται/tai & `future passive' \\
ε/e & & λυ/ly & σε/se & $\varnothing$ & $\varnothing$ & `aorist active' \\
ε/e & & λυ/ly & σα/sa &$\varnothing$ & το/to & `aorist middle' \\
ε/e & & λυ/ly & θη/t\super hε: & & $\varnothing$ & `aorist passive' \\
 & λε/le & λυ/ly & κ/k & ε/e & $\varnothing$ & `perfect active' \\
 & λε/le & λυ/ly &  & & ται/tai & `perfect middle-passive' \\
ε/e & λε/le & λυ/ly & κ/k & ε/e & ι/i & `pluperfect active' \\
ε/e & λε/le & λυ/ly &  & & το/to & `pluperfect middle-passive' \\
 & λε/le & λυ/ly & σ/s & ε/e & ται/tai & `future perfect middle-passive' \\
\end{tabularx}

\normalsize
\subsection{Voice Features}
   
There is also essentially no consensus as to what Voice features should be assigned to middle and passive, regardless of what realizes them. In (8) I give a break down of recent claims on the matter.

\noindent \hfill \par
\pex \textit{Voice feature claims}
\a Embick (1997): Voice is [act] or [nonact, pass]; the middles/passives have Voice features [nonact, passive]
\a Grestenberger (2015, 2019): Voice is [+/--D]; middles/passives have Voice [--D], but the aorist passive has no Voice head / no Voice features
\a Schreiner (2021): Voice can be [+/--active]; middles/passives take [--active]
\xe

\noindent \hfill \par
I'm going to follow Embick in assuming a three-part system for actives, middles and passives, but I will do this using only two features, like Schreiner. Namely I will assume that active voice refers to the feature [+act], middle voice, including so-called middle-passive forms, refers to the feature set [+act, --act], and, wherever there is a separate passive voice, it refers to the feature [--act]. The decision for middle voice is perhaps counter-intuitive, but there is no \textit{a priori} requirement that function composition of synsem features operate over a boolean logic. Even though they are modelled after phonological features, their composition need not continue that analogy. Furthermore, the choice is motivated on at least three grounds. 

First, it provides a natural explanation for the Greek inheritance, in the present and perfect, of the single PIE mediopassive form as distinct from the active form, namely that the middle-passive forms have two features that can be ``activated' according to the syntactic/semantic needs at hand. Second, it provides an easy way for the passive to become separate from the middle (a development in Greek) in the aorist and future systems, namely a separation of the shared feature. Third, and similarly, it provides a way of formally distinguishing the middle and passive in the Greek verb (cf. Alexiadou \& Doron 2012), which is necessary to account for the Greek aorist/future passive forms (assuming that the -θη- [-t\super hε:-] suffix instantiates voice.\footnote{Grestenberger's claim that the -θη- [-t\super hε:-] suffix reflects an absence of Voice is poorly empirically movitated and extremely counter-intuitive.}

Finally, most importantly, it accords well with the empirical facts about the middle versus passive voice's semantic distribution in Attic. Ram\'{o}n (2014) characterizes the aorist passive as a ``restriction'' of the middle. Therein, the -θη- [-t\super hε:-] ending provides a ``passive reading'' as distinct from the ``transitive reading'' of the (active or middle) aorist -σα- [-sa-], and the ``middle -σα- [-sa-] aorists are prototypically agentive and their subject is mostly volitional/agentive''. Similarly, Tronci (2018) shows that the aorist -θη- [-t\super hε:-] transforms the Greek (aorist) verbal system from a bipartite to tripartite system such that the -θη- [-t\super hε:-] aorist passive take sover ``intransitive'' constructions, including passive and unaccusative constructions, while the middle -σα- [-sa-] assumes the role of ``middle-transitive'' constructions, including ``possessive'' and ``benefactive'' constructions. Again, the feature distribution above provides an easy route for this bipartite system to become a tripartite system. It also captures the somewhat odd, mixed semantics of the middle voice distribution, as well as the fact, per Tronci, that deponent verbs (active meanings with middle-passive endings) typically use the middle -σα- [-sa-] forms instead of the passive -θη- [-t\super hε:-] forms.

I assume, as is usual in the above accounts that the aorist realizes the feature [+perfective]. However, Ram\'{o}n (2014) also suggests that the aorist/future passives differ slightly from the middles w/r/t their aspect. Thus, I will also claim that the passive marker realizes the aspect feature [+telic] in addition to [+perfective]. This will be important later for vocabular insertion.

\section{Solution}

Recall the general schema for the Greek verb from (6). The immediate problem to deal with is that Voice features are always realized on the Agr head in some form, even though Voice is structurally well below Agr. There are at least three possible ways to deal with this: i) move Voice or Agr heads, ii) inward-sensitive vocabulary insertion for the Agr head, iii) redundancy rules for Agr that that occur prior to VI (cf. Halle \& Marantz 1993). Option (i) would make little sense: agreement always shows up at the right edge of the verb and must be last, and if the aorist/future passive -θη- [-t\super hε:-] and perfect -κ- [-k-] realize Voice, then Voice has to be lower than Agr. Option (ii) is feasible, but it contradicts the constraint from Bobaljik (2000) that states that inward sensitivity operates over post-insertion vocabulary items, i.e., that inward sensitivity can't ``see'' synsem features.\footnote{Schreiner (2021) argues for and Grestenberger (2015, 2019) assumes sysnem-based inward sensitivity, however.} Thus, option (iii) seems to make the most sense. I provide these redundancy rules in (9). There also needs to be a redundancy rule for the Mood head so that the theme vowel, prior to Agr, only occurs in the proper forms. I give a preliminary version in (10). Note however, that this will fail to get the correct theme/Mood vowel for the perfect system. The perfect system is much more complicated w/r/t to the theme vowel, so I leave it for another work.\footnote{Cf. Schreiner (2021), as well.}

\noindent \hfill \par
\ex \textit{Agr redundancy rules}

\begin{math}
\begin{array}{llll}
    \textrm{a.} & \textrm{Agr} & \rightarrow & \textrm{Agr[+act, --act]} / \underline{\hspace{1em}} \textrm{[+act] [F]* [--act] [F]*} \\
    \textrm{b.} & \textrm{Agr} & \rightarrow & \textrm{Agr[+past]} / \underline{\hspace{1em}} \textrm{[F]* [+past]} \\
\end{array}
\end{math}
\xe


\noindent \hfill \par
\ex \textit{Mood redundancy rules}

\begin{math}
\begin{array}{llll}
    \textrm{a.} & \textrm{Mood} & \rightarrow & \textrm{Mood[--past]} / \underline{\hspace{1em}} \textrm{[--past]} \\
\end{array}
\end{math}
\xe

\noindent \hfill \par
What the redundancy rules in (9) state is that the Agr head gains the features [+act, --act] and/or [+past] in the context of those respective features. I use [F]* with a Kleene star to mean ``any number of any arbitrary feature sets'', i.e, these rules operate over spans of arbitrary length.\footnote{Note that something like this would also be required using option (ii) above.} Theoretically this would include heads that occur after Agr, but I assume this is not possible (assuming, furthermore, that these contexts are only defined for a certain domain of a particular syntactic head). This is because the Agr endings are only sensitive to T and Voice, any other intervening heads don't matter. At this point the first [F]* is unnecessary, since [+act] and [--act], if present, show up on the same Voice head (for the middles and middle-passives), but it will become necessary shortly. These rules account for the cases, like the `present middle-passive' in (7), in which there is only the theme vowel and inflectional endings and no other overt heads, as well as cases, like the `aorist middle' in (7), in which there is some set of intervening heads. The rules apply separately and therefore can both be applied when the features are met (again, like the `aorist middle'). The rule in (10) states that the Mood head gains a [--past] feature in a [--past] context. 

It's possible now to state the Vocabulary items for the Greek regular thematic verb.\footnote{I include only the indicative forms here, though most of the subjunctive and optative forms could be added fairly easily. The future passive optative would present a special challenge insofar it looks like, at least superficially, that the θη t\super hε: suffix is somehow split between Voice and Agr. Furthermore, I only include Agr VIs for the 3rd person forms, for reasons of space, and certain complications/syncretisms therein would still need to be worked out. Finally (11c,d) overgenerate the κ [k] affix for the perfect middle-passive forms, another problem I won't deal with explicitly here. One could however, imagine another impoverishment rule that striles [+act] from the perfect middle-passive forms, or perhaps features of the form [+F, --F] could be considered as distinct from [+F] and [--F], such that [+/--F] is not a subset of [+F, --F]. Both of these options might undermine the motivation for [+F, --F] features in the first place, however.}

\noindent \hfill \par
\ex \textit{vocabulary items}

\noindent \hfill \par
\begin{math}
\begin{array}{l|l|lll}
\textrm{root} & \textrm{a.} & \textrm{λυ [ly]} & \leftrightarrow & \sqrt{\textrm{kill}} \\
\hline 
\textrm{Voice} & \textrm{b.} & \textrm{θη [t\super hε:]} & \leftrightarrow & \textrm{Voice[--act] Asp[+pfctv, +telic]} \\
& \textrm{c.} & \textrm{κα [ka]} & \leftrightarrow & \textrm{Voice[+act]} / \underline{\hspace{1em}} \textrm{Asp[+pfct]} \\
& \textrm{c.} & \textrm{κε [ke]} & \leftrightarrow & \textrm{Voice[+act]} / \underline{\hspace{1em}} \textrm{Asp[+pfct ]Agr[3sg]} \\
\hline 
\textrm{Aspect} & \textrm{d.} & \textrm{σα [sa]} & \leftrightarrow & \textrm{Asp[+pfctv]} / \underline{\hspace{1em}} \textrm{T[+past]}\footnote{I ignore for now the fact that the 3rd person singular active of the aorist and perfect have, respectively -σε [-se] and -κε/[-ke] instead of -σα [-sa] and -κα [-ke].} \\ 
& \textrm{d.} & \textrm{σε [se]} & \leftrightarrow & \textrm{Asp[+pfctv]} / \underline{\hspace{1em}} \textrm{T[+past] Agr[3sg]} \\
& \textrm{e.} & \textrm{RED} & \leftrightarrow & \textrm{Asp[+pfct]} \\
\hline 
\textrm{T} & \textrm{f.} & \textrm{ε [e]} & \leftrightarrow & \textrm{T[+past]} \\
& \textrm{g.} & \textrm{σ [s]} & \leftrightarrow & \textrm{T[fut]} \\
\hline
\textrm{mood} & \textrm{f.} & \textrm{ε [e]} & \leftrightarrow & \textrm{Mood[indic, --past]} \\
& \textrm{h.} & \textrm{ο [o]} & \leftrightarrow & \textrm{Mood[indic, +pl, --2, --past]} \\
\hline
\textrm{Agr} & \textrm{i.} & \textrm{ι [i]} & \leftrightarrow & \textrm{Agr[3sg]} \\
& \textrm{j.} & \textrm{ται [tai]} & \leftrightarrow & \textrm{Agr[3sg, +act, --act]} \\
& \textrm{k.} & \textrm{$\varnothing$} & \leftrightarrow & \textrm{Agr[3sg, +past]} \\
& \textrm{l.} & \textrm{το [to]} & \leftrightarrow & \textrm{Agr[3sg, +act, --act, +past]}\\
\hline
\textrm{elsewhere} & \textrm{l.} & \textrm{$\varnothing$} & \leftrightarrow & \textrm{[\space]}
\end{array}
\end{math}
\xe

\noindent \hfill \par
The VIS operate using the usual Panini's rule: more complex rules beat less complex rules in competition. Note that none of the VIs, though, require inward sensitivity based on synsem features. Still, spans are used in several ways, including the contexts of several VIs (11c,d). Most importantly, the aorist/future passive θη [t\super hε:] (11b) affix itself realizes a span comprising Voice and Aspect. This accounts for the fact that the θη [t\super hε:] and aorist σα [sa] affixes never co-occur, though the former can occur with the future tense. As suggested above, [+telic] with θη/[t\super hε:] is necessary so that it's not inserted in the aorist and future middle forms (based on the subset principle). Furthermore, the VIs here have the result that the future passive has a different aspect realization than the future active and middle. This result is welcome given that data seems to indicate that this distinction holds empirically (Grestenberger 2015, 2019).


I follow Schreiner (2021) in assuming that RED is a reduplication variable that can be filled in according to the root's phonology. The null affix gets inserted for all the non-specified features of the paradigm.\footnote{Namely, [--past], [--pfctv], [+act]} Unfortunately a second null affix was required to account for the 3rd person secondary active Agr ending (11k). The other person/number combinations of the secondary active forms, though, coupled with the redundancy rule in (9b) are not null, however.

Recall, though, the crux of the problem: the future passive gets the otherwise expected middle-passive endings. As it stands, based on the redundancy rule in (9a) and the VIs in (11b,i) the future passive, like the aorist passive, will get the active (primary) endings. I will note, also, that it's very common in Attic for the future paradigm to be deponent (active meaning with passive forms), even when the rest of the verb's paradigm is not deponent. To deal with the future passive facts, as well as the more general future facts, I propose two more redundancy rules.

\noindent \hfill \par
\ex \textit{future passive redundancy rules}

\begin{math}
\begin{array}{llll}
    \textrm{a.} & \textrm{T [+fut]} & \rightarrow & \textrm{T[+fut, +act]} / \underline{\hspace{1em}} \textrm{R$_{ND}$ [--act] [+pfctv, +telic] [fut]} \\
    \textrm{b.} & \textrm{T [+fut]} & \rightarrow & \textrm{T[+fut, +act, --act]} / \underline{\hspace{1em}} \textrm{[fut]} \\
\end{array}
\end{math}
\xe

\noindent \hfill \par
The term R$_{ND}$ here refers to ``non-deponent'' roots, in order to deal with the fact that whether the future is deponent or not is lexically dependent somehow (λύω [ly\textipa{O}:] is such a verb). Thus, (12a) is a redundancy rule that picks out the future passive of certain roots and assigns [+act] to the T node. Assuming that these redundancy rules in (12) occur before those in (9,10), this means that the future passive meets the requirements of (9a), so that Agr heads gets [+act, --act] and the right inflectional endings are inserted. On the other hand the aorist passive still does not meet the requirements of (9a), so it still gets the otherwise odd active endings. The more general case in (12b) states that the future tense is flatly assigned [+act, --act], i.e., middle, features. This means the future being deponent is the general case, and (12a) and (12b) compete under Panini's principle, just like VIs do. Again, that the deponent case is the unmarked case accords well with the general trend in Attic for futures to be deponent. The rules in (12) also make it easy to account for the difference between Attic and Doric dialects. Namely, the Doric future passive takes the active endings just like the aorist passive, and deponent futures are rarer (Grestenberger 2015, 2019). This means that Doric simply doesn't have these redundancy rules in (12). The structure for the 3rd person singular future passive of λύω [ly\textipa{O}:] is shown in (13).

\ex \textit{future passive}
    
    \footnotesize
    \begin{forest}
    [Agr
    [Mood
    [Tense
    [Aspect
    [Voice
    [V [λυ]]
    [Voice,
    tikz={
            \node [inner sep=0pt, fit to=tree] (D3) {};
            \draw [decoration={brace, amplitude=6pt, mirror},
                decorate, semithick]
                (D3.south west) --++ (4.5,0)
            node [xshift=-5.4em, yshift=-1.5em] {θη};
            }
    [{[--act]}]]]
    [Aspect [{[+pfctv, +telic]}]]]
    [T[{[--pres, fut, +act]} \\ σ]]]
    [Mood[{[indic, --pres]} \\ ε]]]
    [Agr[{[3sg, +act, --act]} \\ ται]]]
    \end{forest}
    
    \xe


In any case, the VIs in (11) and the redundancy rules in (9,10,12) can account reasonably well for the strange aorist/future passive facts in Attic Greek. A number of details still need to be worked out, most notably the precise Agr inflectional ending VIs and mood/theme vowel VIs. The general idea, though, that span-based redundancy rules and VI underspecification, without recourse to inward sensitivy to synsem features, can account for these facts, seems to hold.

\section{Conclusion}

This paper has presented a DM approach to a curious puzzle in Attic Greek: the fact that aorist passive forms take active endings while the related future passive takes middle-passive endings. The paper uses underspecification of vocabulary items and redundancy rules, both of which use spans and require Panini's principle, to account for this strange data. That redundancy rules might operate using the same principle, Panini's principle, as VIs do is an interesting and fruitful result, and represents an advantage over an analysis that chooses inward sensitivity to synsem features over redundancy rules: redundancy rules still follow a more general principle while allowing cyclic vocabulary insertion (another general feature of the system) to remain in place. Furthermore, redundancy rules provide a natural way of producing diachronic change that's not ``analogy'': these rules could hypothetically be added relatively easily to the morphology component of the grammar. This is especially true if the redundancy rules utilize language-specific innovations, like the distinct passive in aorist Greek. If the [+act, --act] account is on the right track, the passive could split off fairly easily diachronically, and redundancy rules could then target it easily. Similarly, the prevalence of Attic future deponents can be explained by this rule being added in the Morphology. Furthermore, these rules, unlike dissociated morphemes, don't require the morphology to be able to change the structural output of the syntax, only to change featural content, a much less powerful claim. Of course, the possibility that the feature structure [+F, --F] is possible needs to be further explored, as it has wide-ranging consequences on the morphological component and how synsem features are organized. Finally, as mentioned, a number of the precise details about the Attic inflectional system still need to be worked out. In any case, this seems like a fruitful first step for this longstanding puzzle.

\small
\nocite{*}
\printbibliography

\end{document}